\scanpage{201}
$\mathbb{C}$, i.e. $\overline{k(a)}\simeq\mathbb{R}$\inserted{isomor. can.}, alors $\unsure{\xi|\zeta}_{x,\mathbb{Z}}$ est encore défini par la valuation archimédienne sur $k(a)$, mais \deleted{suivant que les} \inserted{\unsure{conformt aux|conformément aux|compte tenu}} choix possibles dans la clôture algébrique de $k(a)$, le gpe de Galois de $\mathbb{C}=\overline{\mathbb{R}}$ sur $\mathbb{R}$ opère dans $\xi_{x,\mathbb{Z}}$, soit \uline{$\tau$ cette opération}, qui est involutive. [On ne la confondra pas avec l’\unsure{endom|endom.|action} \uline{$\tfrac12$ linéaire} de $\xi_{\bar{a},\infty}$ correspondant à $\tau$ ; il n’y est pas question par exemple d’un réseau stable].

\margin{On écrit aussi $\xi_{a,\mathbb{Z}}$ au lieu de $\xi_{x,\mathbb{Z}}$,}

% slanted
\margin{on choisit le \uline{\unsure{plon|plus|plgt}} défini par le pt à val dans $\mathbb{C}$}

Dans ce cas, les invariants sont

\diagram{\begin{tikzcd}[ampersand replacement=\&,column sep=small,row sep=small]
\xi_{a,\infty}\in T_{\alpha}(\mathbb{R}) \arrow[dr,dash] \& {}\\
\deleted{\xi_{\unsure{a|\bar a},\infty}\otimes_{\mathbb{R}}\mathbb{C}\simeq} \& \xi_{\bar{a},\infty}\in T_{\alpha}(\mathbb{C})\\
\xi_{\bar{a},\mathbb{Z}}\in T_{\alpha}(\mathbb{Z}) \arrow[ur,dash] \& {}
\end{tikzcd}}

\[
\xi_{\bar{a},\infty}\overset{\text{\uuline{can}}}{\simeq}\xi_{a,\infty}\otimes_{\mathbb{R}}\mathbb{C}\simeq\xi_{\bar{a},\mathbb{Z}}\otimes_{\mathbb{Z}}\mathbb{C}
\]

$\unsure{\Pi|\pi}_1(X)$ n’opère pas à priori dans $\xi_{a,\mathbb{Z}}$, les seules opérations sont celles de $\tau$ si $a$ est \uline{réel}. Cependant, sur $\xi_{\bar{a},\infty}$, il y a (en plus éventuellement de $\tau$, si $a$ est réel) les opérations de $\operatorname{Aut}_{k(a)}^{\unsure{\ell|i}}(k)$, qui sont $k$ linéaires.

% archive sheet number: 200
